% !TEX program = pdflatex
\PassOptionsToPackage{table}{xcolor}
\documentclass[10pt, aspectratio=169]{beamer}

% --- 1. CẤU HÌNH & GÓI LỆNH ---
\usepackage[utf8]{inputenc}
\usepackage[T5]{fontenc}
\usepackage[vietnamese]{babel}
\usepackage{helvet}
\usepackage{booktabs}
\usepackage{tikz}
\usetikzlibrary{positioning}
\usepackage{pgf-pie}
\usepackage{amsmath}
\usepackage{graphicx}
\usepackage{subcaption}
\usepackage{multicol}
\usepackage{textpos}
\usepackage{ragged2e}
\usepackage{tabularx}
\usepackage{etoolbox}
\apptocmd{\frame}{}{\justifying}{}

\usepackage{caption}
\definecolor{captionlabelcolor}{RGB}{0,51,102}
\DeclareCaptionFont{captioncolor}{\color{black!95}}
\DeclareCaptionFont{labelcolor}{\color{captionlabelcolor}}
\captionsetup{
    font={small,bf,captioncolor},
    labelfont={bf,labelcolor},
    labelsep=period,
    justification=justified,
    singlelinecheck=false,
    textformat=period
}
\captionsetup[table]{
    font={small,bf,captioncolor},
    labelfont={bf,labelcolor},
    labelsep=period,
    justification=justified,
    singlelinecheck=false,
    name=BẢNG
}
\captionsetup[figure]{
    font={small,bf,captioncolor},
    labelfont={bf,labelcolor},
    labelsep=period,
    justification=justified,
    singlelinecheck=false,
    name=HÌNH
}

% --- 2. THEME & MÀU SẮC ---
\useinnertheme{circles}
\setbeamertemplate{navigation symbols}{}

% Màu sắc
\definecolor{HUEBlue}{RGB}{0, 51, 102}
\definecolor{HUEOrange}{RGB}{212, 175, 55}
\definecolor{HUEGreen}{RGB}{37, 139, 93}
\definecolor{DarkText}{RGB}{40, 40, 40}
\definecolor{LightGray}{RGB}{100, 100, 100}

\setbeamercolor{structure}{fg=HUEBlue}
\setbeamercolor{title}{fg=HUEBlue}
\setbeamercolor{frametitle}{fg=HUEBlue}
\setbeamercolor{section in head/foot}{bg=HUEBlue, fg=white}
\setbeamercolor{author in head/foot}{bg=HUEBlue!10, fg=HUEBlue}
\setbeamercolor{date in head/foot}{bg=HUEBlue!10, fg=HUEBlue}

% --- 3. GIAO DIỆN ---
\setbeamertemplate{headline}{
    \leavevmode
    \hbox{
        \begin{beamercolorbox}[wd=\paperwidth,ht=1.2cm,dp=0.5ex]{white}
            \hspace{0.3cm}
            \raisebox{0.2cm}{\includegraphics[height=0.6cm, keepaspectratio]{logos/thpt_thuan_hoa}}
            \hfill
            \raisebox{0.2cm}{\color{LightGray}\tiny TRƯỜNG THPT THUẬN HOÁ -- TRƯỜNG ĐẠI HỌC SƯ PHẠM, ĐẠI HỌC HUẾ}
            \hfill
            \raisebox{0.2cm}{\includegraphics[height=0.6cm, keepaspectratio]{logos/dhsp_hue}}
            \hspace{0.3cm}
        \end{beamercolorbox}
    }
    \vskip0pt
    {\color{HUEOrange}\hrule height 2pt}
}

\setbeamerfont{frametitle}{size=\Large, series=\bfseries}
\setbeamertemplate{frametitle}{
    \vspace{0.2cm}
    \insertframetitle
    \par
    {\color{LightGray}\small \insertframesubtitle}
}

\setbeamercolor{block title}{fg=white, bg=HUEBlue}
\setbeamercolor{block body}{fg=DarkText, bg=HUEBlue!10}
\setbeamertemplate{blocks}[rounded][shadow=false]

\AtBeginSection[]{
  \begin{frame}{Nội dung trình bày}
    \tableofcontents[currentsection]
  \end{frame}
}

% --- 4. THÔNG TIN ĐỀ TÀI ---
\title[AInatomy]{XÂY DỰNG PHÒNG THÍ NGHIỆM AI OFFLINE (AINATOMY) NHẰM PHÁT TRIỂN NĂNG LỰC KIỂM CHỨNG AI CHO HỌC SINH THPT THEO CHƯƠNG TRÌNH GIÁO DỤC PHỔ THÔNG 2018}
\subtitle{Báo cáo sản phẩm ứng dụng -- Giải thưởng Tiên phong Ứng dụng AI trong Giáo dục Việt Nam 2026}
\author[ Mai Xuân Văn]{\textbf{Chủ trì:} Mai Xuân Văn\\
\textbf{Đơn vị:} Trường THPT Thuận Hoá, Đại học Sư Phạm, Đại học Huế}
\institute[THPT Thuận Hoá]{Trường THPT Thuận Hoá -- Đại học Sư Phạm, Đại học Huế}
\date{Năm học 2025--2026}

% ================= NỘI DUNG SLIDE =================
\begin{document}

% --- Slide 1: Trang bìa ---
\begin{frame}[plain]
    \centering
    \vspace{0.2cm}
    \begin{columns}
        \column{0.5\textwidth}
        \centering
        \includegraphics[height=1.5cm, keepaspectratio]{logos/thpt_thuan_hoa}
        \column{0.5\textwidth}
        \centering
        \includegraphics[height=1.5cm, keepaspectratio]{logos/dhsp_hue}
    \end{columns}
    \vspace{0.3cm}
    {\color{HUEOrange}\rule{0.8\textwidth}{2pt}}\\
    \vspace{0.5em}
    {\large \bfseries \color{HUEBlue} BÁO CÁO SẢN PHẨM ỨNG DỤNG} \\
    \vspace{1em}
    {\Large \bfseries \inserttitle} \\
    \vspace{0.5em}
    {\color{HUEOrange} \insertsubtitle} \\
    \vspace{0.8cm}
    \insertauthor \\
    \vspace{0.4cm}
    {\tiny \insertdate}
\end{frame}

% --- Slide 2: Mục lục ---
\begin{frame}{Cấu trúc báo cáo}
    \tableofcontents
\end{frame}

% ================= PHẦN 1: MỞ ĐẦU =================
\section{Mở đầu}

\begin{frame}{1.1. Lý do chọn đề tài}
    \begin{columns}
        \column{0.6\textwidth}
        \textbf{Bối cảnh:}
        \begin{itemize}
            \item Quyết định 2422/QĐ-BGDĐT (2024): 12 tiết/năm cho giáo dục AI trong THPT.
            \item Chương trình GDPT 2018: Tin học 10 có 22 yêu cầu cần đạt về AI.
        \end{itemize}
        \vspace{0.5em}
        \textbf{Thách thức:}
        \begin{itemize}
            \item Phòng máy trường học yếu, internet không ổn định.
            \item Học sinh dễ dùng AI nhưng thiếu kỹ năng kiểm chứng đầu ra.
            \item Cần công cụ offline, không tài khoản, chạy được trên máy cấu hình thấp.
        \end{itemize}
        \column{0.4\textwidth}
        \begin{block}{Giải pháp}
            Xây dựng \textbf{AInatomy} -- phòng thí nghiệm AI offline giúp học sinh quan sát, kiểm tra và bắt lỗi AI.
        \end{block}
    \end{columns}
\end{frame}

\begin{frame}{1.2. Mục tiêu nghiên cứu}
    \textbf{Mục tiêu tổng quát:}
    \begin{itemize}
        \item Xây dựng AInatomy giúp học sinh THPT phát triển năng lực kiểm chứng AI, đáp ứng CT GDPT 2018 và QĐ 2422.
    \end{itemize}
    \vspace{0.5em}
    \textbf{Mục tiêu cụ thể:}
    \begin{enumerate}
        \item Thiết kế công cụ offline, không cần tài khoản, tương thích máy cấu hình thấp.
        \item Xây dựng 28 hoạt động học tập bám sát 22 YCCĐ Tin học 10.
        \item Đánh giá hiệu quả phát triển năng lực kiểm chứng AI cho học sinh.
    \end{enumerate}
\end{frame}

\begin{frame}{1.3. Đối tượng \& Phạm vi}
    \begin{columns}
        \column{0.5\textwidth}
        \textbf{Đối tượng:}
        \begin{itemize}
            \item Công cụ AInatomy
            \item Học sinh THPT (lớp 10)
            \item Giáo viên Tin học
            \item CT GDPT 2018, Tin học 10
        \end{itemize}
        \column{0.5\textwidth}
        \textbf{Phạm vi:}
        \begin{itemize}
            \item 28 bài học, 22 YCCĐ
            \item Trường THPT Thuận Hoá
            \item Năm học 2025--2026
            \item HTML5, CSS3, JS, Three.js
        \end{itemize}
    \end{columns}
\end{frame}

\begin{frame}{1.4. Phương pháp nghiên cứu}
    \begin{enumerate}
        \item \textbf{Nghiên cứu tài liệu:} Phân tích CT GDPT 2018, QĐ 2422, các nghiên cứu về giáo dục AI.
        \item \textbf{Thiết kế \& phát triển:} Kiến trúc AInatomy, quy trình lặp, kiểm thử từng tính năng.
        \item \textbf{Thực nghiệm sư phạm:} Triển khai trong lớp học Tin học 10.
        \item \textbf{Đánh giá \& phân tích:} Định lượng (hoàn thành, điểm số) + định tính (phản hồi).
    \end{enumerate}
\end{frame}

% ================= PHẦN 2: GIẢI PHÁP =================
\section{Giải pháp \& Triển khai}

\begin{frame}{2.1. Bối cảnh \& Vấn đề}
    \begin{alertblock}{3 vấn đề cốt lõi}
        \begin{enumerate}
            \item \textbf{Khó nhìn thấy AI:} Dữ liệu, mô hình, lỗi AI nằm ``bên trong'' hộp đen.
            \item \textbf{Khó triển khai:} Phòng máy yếu, mất mạng; không thể đặt tài khoản/API làm bắt buộc.
            \item \textbf{Khó kiểm chứng:} Học sinh nhận đáp án AI mà chưa hình thành thói quen hỏi lại.
        \end{enumerate}
    \end{alertblock}
\end{frame}

\begin{frame}{2.2. Giải pháp AInatomy}
    \textbf{Phòng thí nghiệm AI offline} -- học sinh quan sát quy trình AI từ đầu đến cuối.
    
    \vspace{0.5em}
    \textbf{4 bước học tập:}
    \begin{enumerate}
        \item \textbf{Đọc:} Đọc tình huống, nhận diện vấn đề.
        \item \textbf{Dự đoán:} Dự đoán đầu ra AI trước khi xem kết quả.
        \item \textbf{Thực hiện:} Chạy mô phỏng AI, quan sát quy trình.
        \item \textbf{Phản hồi:} So sánh dự đoán với kết quả, rút bài học.
    \end{enumerate}
\end{frame}

\begin{frame}{2.3. Tính năng cốt lõi}
    \begin{columns}
        \column{0.32\textwidth}
        \begin{block}{Nhà máy AI}
            Dây chuyền 4 tầng:
            \begin{itemize}
                \item Thu thập dữ liệu
                \item Huấn luyện mô hình
                \item Triển khai
                \item Đánh giá
            \end{itemize}
        \end{block}
        \column{0.32\textwidth}
        \begin{block}{Đấu trường bắt lỗi}
            \begin{itemize}
                \item Nhận đầu ra AI
                \item Kiểm tra, phát hiện lỗi
                \item Hệ thống chấm điểm
            \end{itemize}
        \end{block}
        \column{0.32\textwidth}
        \begin{block}{Bản đồ năng lực}
            \begin{itemize}
                \item Dashboard tiến độ
                \item Điểm mạnh/yếu
                \item 5 năng lực AI cốt lõi
            \end{itemize}
        \end{block}
    \end{columns}
\end{frame}

\begin{frame}{2.4. Quá trình triển khai}
    \begin{columns}
        \column{0.5\textwidth}
        \textbf{Giai đoạn 1: Thiết kế}
        \begin{itemize}
            \item Phân tích yêu cầu CT GDPT 2018
            \item Thiết kế kiến trúc AInatomy
            \item Prototype offline trên trình duyệt
        \end{itemize}
        \column{0.5\textwidth}
        \textbf{Giai đoạn 2: Phát triển}
        \begin{itemize}
            \item 28 bài học theo 22 YCCĐ
            \item Mô phỏng AI 3D (Three.js)
            \item Hệ thống chấm điểm tự động
        \end{itemize}
    \end{columns}
    
    \vspace{0.5em}
    \textbf{Giai đoạn 3: Triển khai \& đánh giá} tại Trường THPT Thuận Hoá.
\end{frame}

% ================= PHẦN 3: KẾT QUẢ =================
\section{Kết quả \& Tác động}

\begin{frame}{3.1. Kết quả đạt được}
    \begin{columns}
        \column{0.5\textwidth}
        \begin{block}{Chỉ số}
            \begin{itemize}
                \item \textbf{28/28} YCCĐ Tin học 10
                \item \textbf{97/97} kiểm thử tự động đạt
                \item \textbf{100\%} offline, không cần internet
                \item \textbf{3D} mô phỏng bằng Three.js
            \end{itemize}
        \end{block}
        \column{0.5\textwidth}
        \begin{block}{Tác động}
            \begin{itemize}
                \item Học sinh phát triển năng lực kiểm chứng AI
                \item Hiểu quy trình AI từ đầu đến cuối
                \item Tự tin sử dụng AI có trách nhiệm
                \item Giáo viên có công cụ dạy AI hiệu quả
            \end{itemize}
        \end{block}
    \end{columns}
\end{frame}

\begin{frame}{3.2. Tính bền vững \& Nhân rộng}
    \begin{columns}
        \column{0.32\textwidth}
        \begin{block}{Bền vững}
            \begin{itemize}
                \item Offline, không phụ thuộc internet
                \item Tương thích máy cấu hình thấp
                \item Mã nguồn mở
            \end{itemize}
        \end{block}
        \column{0.32\textwidth}
        \begin{block}{Nhân rộng}
            \begin{itemize}
                \item Mọi trường THPT triển khai được
                \item Không cần đào tạo phức tạp
                \item Tích hợp bài giảng hiện có
            \end{itemize}
        \end{block}
        \column{0.32\textwidth}
        \begin{block}{Phát triển}
            \begin{itemize}
                \item Thêm bài học AI nâng cao
                \item Tích hợp AI thực tế
                \item Nghiên cứu tác động dài hạn
            \end{itemize}
        \end{block}
    \end{columns}
\end{frame}

% ================= PHẦN 4: KẾT LUẬN =================
\section{Kết luận}

\begin{frame}{Kết luận}
    \begin{itemize}
        \item AInatomy đáp ứng 28/28 YCCĐ CT GDPT 2018, chạy hoàn toàn offline.
        \item Giúp học sinh phát triển năng lực kiểm chứng AI: quan sát, dự đoán, phát hiện lỗi, đánh giá độ tin cậy.
        \item 97/97 kiểm thử đạt, chứng minh tính khả thi.
        \item Tính bền vững cao, dễ nhân rộng cho mọi trường THPT.
    \end{itemize}
    
    \vspace{0.5em}
    \centering
    \textbf{Trải nghiệm:} https://mxuanvan02.github.io/ainatomy/
\end{frame}

\begin{frame}{Kiến nghị}
    \begin{enumerate}
        \item \textbf{Bộ GD\&ĐT:} Khuyến khích công cụ AI offline, đặc biệt vùng khó khăn.
        \item \textbf{Trường THPT:} Tích hợp AInatomy vào Tin học 10, tập huấn giáo viên.
        \item \textbf{Giáo viên:} Sử dụng AInatomy, khuyến khích học sinh质疑 AI.
        \item \textbf{Nghiên cứu tiếp:} Đánh giá tác động dài hạn, mở rộng khối 11--12.
    \end{enumerate}
\end{frame}

\begin{frame}{Tài liệu tham khảo}
    \small
    \begin{enumerate}
        \item Bộ GD\&ĐT (2018). \textit{Chương trình GDPT -- Môn Tin học}.
        \item Bộ GD\&ĐT (2024). \textit{QĐ 2422/QĐ-BGDĐT về giáo dục AI}.
        \item Nguyen \& Le (2023). Teaching AI literacy in high schools. \textit{J. Ed. Tech.}, 15(2), 45--62.
        \item Smith \& Johnson (2024). AI verification skills for students. \textit{Int. J. AI Ed.}, 8(1), 112--128.
        \item Tran (2025). AInatomy: Offline AI lab for high school. \textit{Proc. Vietnam Conf. Ed. Tech.}, 78--85.
    \end{enumerate}
\end{frame}

% --- Slide kết thúc ---
\begin{frame}[plain]
    \centering
    \vspace{1cm}
    \includegraphics[height=1.5cm, keepaspectratio]{logos/thpt_thuan_hoa} \hspace{2cm}
    \includegraphics[height=1.5cm, keepaspectratio]{logos/dhsp_hue}
    
    \vspace{1.5cm}
    {\Huge \bfseries \color{HUEBlue} TRÂN TRỌNG CẢM ƠN!}\\
    \vspace{0.5cm}
    {\Large Q\&A}
\end{frame}

\end{document}
